% =====================================================================
\part{A banca}
% =====================================================================

Cinquenta perguntas, agrupadas por origem provável. Responda cada uma em voz
alta antes de ler a resposta sugerida --- o valor deste capítulo está no
exercício, e não na leitura.

As respostas estão escritas em primeira pessoa e no registro em que devem ser
ditas: direto, sem hedge desnecessário, e com o limite declarado \emph{dentro}
da resposta e não como uma desculpa no fim.

\chapter{Sobre o dado e a medida}

\pergunta{1. O MapBiomas é confiável? Qual é a acurácia da Coleção 10.1 em
Goiás?}
\resposta{Reconheça o limite sem se defender: a acurácia do classificador é a
única fonte de erro que \emph{nenhum} recorte deste trabalho alcança --- está
declarado no quadro de graus de estabelecimento como o resíduo da Frente 1.
Diga em seguida o que você fez a respeito: toda afirmação-manchete que passa
por classe do satélite é conferida contra uma fonte que não passa pelo
classificador (área plantada da PAM, rebanho da PPM, crédito do SICOR), e a
supressão de vegetação foi validada externamente contra o PRODES. Foi
justamente essa disciplina que revelou a migração de rótulo do Mosaico.}

\pergunta{2. Por que a Coleção 10.1, e não a mais recente?}
\resposta{A dissertação está fechada na 10.1, que é a coleção sobre a qual todo
o encadeamento de auditorias foi construído. Trocar de coleção no meio implica
refazer o censo de pixels, a matriz de transição, as decisões D25--D26 e todas
as verificações que dependem delas. É trabalho de outra rodada, e não um ajuste.}

\pergunta{3. O que exatamente é a classe ``Mosaico de Usos'', e por que ela
virou o problema central?}
\resposta{É a classe em que o classificador não separou usos --- ela contém
sistemas integrados de lavoura e pecuária, mosaicos finos demais para separar a
30 metros, e mosaico antigo e estável. Ela virou central porque, no fim da
série, parte da conversão de pastagem em lavoura passou a ser rotulada como
Mosaico em vez de Agricultura. Isso faz a conversão parecer cair 88\% no Sul
enquanto a soja que o IBGE mede em campo cresce 38\%.}

\pergunta{4. Por que você não simplesmente somou Agricultura e Mosaico?}
\resposta{Porque somá-las suporia que todo o Mosaico é agricultura mal
rotulada, o que é falso. A união superconta tanto quanto a classe isolada
subconta, e responde a uma pergunta diferente e mais grossa: ``quanta terra saiu
de pastagem para lavoura ou uso misto?''. Por isso reporto o intervalo entre as
duas réguas e considero robusta apenas a conclusão que sobrevive nos dois
extremos.}

\pergunta{5. Como você sabe que a migração de rótulo não está inflando os seus
resultados?}
\resposta{Porque ela aponta na direção contrária. Entre 2019 e 2024, quatro
medidas indicam deslocamento ao norte e a quinta indica que nada se moveu ---
justamente a única cujo rótulo mudou. E o centro da massa reetiquetada fica 46
km ao \emph{norte} do centro da agricultura visível, com crescimento por região
correlacionado 0,84 com a soja do IBGE. A régua crua faz a marcha parecer mais
fraca do que foi.}

\pergunta{6. E o fluxo reverso, de pastagem para vegetação natural? Não é
regeneração?}
\resposta{Não, e testei. De 75\% a 98\% dele, conforme o par de anos, dirige-se
à formação savânica, e ele é aproximadamente simétrico ao fluxo de ida --- a
razão entre ida e volta colapsa de 8,4 vezes no primeiro par decenal para 1,22
no par anual mais recente. Regeneração real é lenta e unidirecional; oscilação
de classificador é bidirecional e balanceada. Há um componente real,
minoritário, restrito à formação florestal, que só se acumula em horizontes
longos.}

\pergunta{7. Qual é a resolução da camada de aptidão da Embrapa, e ela é
adequada?}
\resposta{É 1:500.000, com 8.284 polígonos em Goiás, cobrindo 33,0 Mha ou 97\%
do estado. É escala grosseira diante dos 30 metros da cobertura, e por isso ela
resolve \emph{gradiente regional} e não contraste local, e é exatamente como
ela é usada, como exposição por AMC num desenho de interação. Há ainda uma
suposição de medida que declaro: converter o grupo de aptidão em escore por
$7 - \text{grupo}$ trata como cardinal o que na origem é ordinal.}

\pergunta{8. Por que a lotação está em cabeças por hectare e não em Unidade
Animal?}
\resposta{Porque o fator de conversão é constante, e portanto se cancela em
qualquer análise de tendência ou de variação, que é onde a lotação entra aqui.
E por ser fixo no tempo, ele produziria um nível não comparável entre 1985 e
2024, já que o peso médio do animal mudou no período. Reporto a medida
diretamente observada.}

\pergunta{9. Por que as variáveis monetárias anteriores ao Plano Real foram
excluídas em vez de convertidas?}
\resposta{A conversão entre Cruzeiro, Cruzado, Cruzado Novo e Cruzeiro Real
introduziria erro de medida em séries usadas como regressores. Por isso o valor
da produção da PAM e a série de leite em valor não entram no painel, embora as
quantidades físicas correspondentes entrem. Todo valor em reais é deflacionado
para reais de dezembro de 2024 pelo IPCA acumulado.}

\pergunta{10. Como você mede área, se a projeção que usa não é equivalente?}
\resposta{Não meço área na projeção. A policônica (EPSG:5880) é de compromisso e
serve à geometria: centroides, distâncias, deslocamentos, recortes. As áreas vêm
da contagem de células de 30 m na grade nativa do produto, convertida em
hectares pelo cosseno da latitude média observada de cada município, sem passar
por projeção nenhuma.}

\chapter{Sobre as unidades e o desenho}

\pergunta{11. Por que 166 AMCs e não os 246 municípios?}
\resposta{Porque 62 dos 246 municípios só aparecem na estatística depois de
1985, e as duas fontes se comportam de maneiras opostas diante de uma
emancipação: o raster já vem no polígono de hoje em todos os anos, e a
estatística do IBGE vem no município de então. Municípios-pai registram quedas
de 50\% a 80\% no rebanho no ano da emancipação, enquanto o total estadual
atravessa esses anos sem descontinuidade. As ondas de 1993 e 1997 caem dentro
das janelas longitudinais e dominariam qualquer estimador de variação.}

\pergunta{12. Qual é o custo dessa agregação?}
\resposta{Os 133 municípios reunidos nos 53 grupos deixam de ser analisáveis
isoladamente nas séries longas, e a malha perde 80 unidades. É um preço
inevitável quando o dado vem de pesquisa: a informação histórica do
município-filho não existe separada, porque estava no pai.}

\pergunta{13. Você usa 166 e 246 em lugares diferentes. Isso não é
inconsistente?}
\resposta{É uma estratégia declarada de dois trilhos, e usar o trilho errado
seria erro metodológico. As 166 AMCs são a malha de toda análise longitudinal;
a malha municipal atual fica reservada às análises transversais e ao período
recente --- o Censo de 2017, os mapas de um único ano, a conta do IFDM. Cada
resultado declara em qual malha foi medido.}

\pergunta{14. Por que cinco mesorregiões em alguns lugares e três regiões em
outros?}
\resposta{A malha de cinco serve à leitura grossa do eixo Sul--Norte, ordenada
pela latitude média dos pixels de conversão. A de três é usada onde o número de
unidades por bloco precisa ser grande o bastante para sustentar comparação: a
decomposição da fronteira, o desenho de \emph{drive} comum e a conta de carbono.
Como ``Sul'' nomeia a mesma coisa nas duas partições mas ``Centro'' e ``Norte''
não, cada resultado declara a malha em que foi medido, e é por isso que os
49\% e os 37\% de queda não são o mesmo número.}

\pergunta{15. Por que não usar as Regiões Geográficas Imediatas, que são 22?}
\resposta{Porque elas não acrescentariam leitura a nenhuma das duas escalas que
já uso. O recorte regional serve à leitura grossa do eixo, para a qual cinco
unidades ordenadas em latitude bastam, e toda a análise fina roda sobre as 166
AMCs, que são mais numerosas que as Imediatas e têm a vantagem de território
constante.}

\pergunta{16. Por que primeira diferença em tudo?}
\resposta{Por duas razões independentes. A primeira é que quase toda série do
painel tem tendência, e correlacionar níveis entre séries com tendência comum
produz associações altas e vazias. A segunda é que muitas variáveis são
relacionadas por construção --- área plantada e produção, ou área e valor
adicionado, com correlação em nível da ordem de 0,9. Auditei isso: as mesmas
duplas têm correlação próxima de zero em variação \emph{within}.}

\pergunta{17. Há exceções à primeira diferença?}
\resposta{Três, cada uma por exigência do método. A correlação em nível estadual
é descritiva e vem com os limites declarados. O Toda-Yamamoto ajusta em níveis
por construção, e é isso que torna a inferência válida sob integração. E as
equações do teto de oferta usam efeitos fixos em vez de diferença, porque o que
elas medem é o que acontece a uma unidade \emph{enquanto ela se esgota} ---
diferenciar apagaria justamente esse nível.}

\chapter{Sobre a estatística e a inferência}

\pergunta{18. O que exatamente os efeitos fixos controlam aqui?}
\resposta{Como a equação está em primeira diferença, eles absorvem coisa
distinta da que absorveriam em nível. A diferenciação já retirou o patamar de
cada município, então o efeito fixo de unidade passa a retirar a
\emph{tendência própria} da unidade. O de ano absorve o choque comum às
variações daquele ano. O que resta para o coeficiente estimar é o quanto a
variação de um município num ano se afasta do que a tendência dele e o ano em
curso já explicavam.}

\pergunta{19. Então os seus coeficientes são causais?}
\resposta{Não, e os chamo de associações bem controladas. O painel controla o
que é fixo no lugar e o que é comum ao ano, mas não um determinante que varie
\emph{dentro} do município ao longo do tempo e tenha ficado fora do modelo.}

\pergunta{20. Por que não um desenho de diferenças em diferenças para os marcos
de política?}
\resposta{Eu rodei, e ele passa nos testes de robustez usuais: tendências
paralelas, estudo de evento, placebo temporal. Mesmo assim eu o rebaixei, porque
robustez não conserta um problema de desenho: os quatro marcos testados são
federais ou de alcance nacional, e alcançaram no mesmo momento o estado tratado
e os controles do mesmo bioma. Sem grupo não tratado, o coeficiente mede
exposição diferencial ao mesmo choque. Separar as hipóteses exigiria variação
dentro do estado, e esse desenho eu não tenho.}

\pergunta{21. As unidades vizinhas não são independentes. Como isso afeta a sua
inferência?}
\resposta{Afeta em dois lugares, e tratei os dois. No painel, o erro-padrão
agrupado por unidade corrige a correlação dentro da unidade, mas não entre
vizinhas no mesmo ano --- então reestimei em SAR e SEM: os parâmetros espaciais
ficam entre $+0{,}35$ e $+0{,}56$, sete das oito estimativas se atenuam no
máximo 14,6\% e uma cresce 0,6\%, nenhuma troca de sinal nem perde
significância. Nos centros de massa, refiz o \emph{bootstrap} por blocos
espaciais, varrendo o tamanho de 1 a 14 AMCs: o intervalo alarga como esperado e
o veredito não muda em nenhum tamanho.}

\pergunta{22. Por que reportar $p$ de permutação em vez do $p$ agrupado?}
\resposta{Porque num desenho \emph{shift-share} com um único \emph{shifter} o
erro-padrão agrupado trata cada ano como realização independente do choque,
quando há um só choque observado 38 vezes e serialmente correlacionado. Isso o
torna otimista. A permutação circular rotaciona a série preservando a
autocorrelação e é a régua defensável. E eu reporto as duas lado a lado, porque
mostrá-las juntas é o que impede a régua de ser trocada em silêncio.}

\pergunta{23. Você diz que o menor $p$ possível é 0,026. Explique.}
\resposta{Uso todas as $T-1$ rotações possíveis, e o $p$ inclui a estatística
observada entre as realizações: é o número de rotações ao menos tão extremas
quanto a observada, mais um, dividido por $T$. Com 38 anos, o menor valor
atingível é $1/38 \approx 0{,}026$. Quando a latitude devolve exatamente isso,
significa que \emph{nenhuma} rotação superou o observado --- é o melhor desfecho
possível do teste, e não uma margem folgada. A resolução do teste é parte do que
se pode afirmar.}

\pergunta{24. Um $R^2$ \emph{within} de 0,001 não é irrelevante?}
\resposta{Ele é pequeno, e eu o reporto. Mas o desenho não se propõe a explicar
a variação do rebanho: com efeito fixo de ano, o nível do choque é absorvido por
construção, e o que resta identificado é apenas o \emph{gradiente} do efeito.
Um termo de interação puro sobre uma variação anual ruidosa não teria como
explicar muito. É também parte da razão de eu classificar essa frente como
corroborante e não estabelecida.}

\pergunta{25. Por que Toda-Yamamoto e não Granger direto?}
\resposta{Porque um teste de Granger aplicado a séries integradas fabrica
precedência inexistente. A pastagem do Norte não é estacionária nem em nível nem
em primeira diferença pelo ADF; só na segunda diferença os dois testes voltam a
concordar, e é essa linha que fixa $d_{\max} = 2$. A agricultura do Sul é
estacionária já em nível. Séries de ordens diferentes não podem ser
cointegradas, e qualquer precedência achada nessa configuração é artefato.}

\pergunta{26. E o que Toda-Yamamoto faz de diferente?}
\resposta{Ajusta o vetor autorregressivo em \emph{níveis} com $p + d_{\max}$
defasagens e aplica o teste de Wald apenas às $p$ primeiras. As defasagens
extras absorvem a não estacionariedade, e ao restringir o teste às primeiras a
estatística recupera a distribuição usual mesmo sob integração.}

\pergunta{27. O seu apêndice mostra um $p < 0{,}001$ no sentido Norte para Sul.
Isso não contradiz a sua tese?}
\resposta{Aquele bloco está na tabela de propósito: foi ele que originou a
decisão D16, e o reporto para que se veja o que o diagnóstico teve de descartar.
As três colunas de $p$ não são a mesma régua. O $p$ de Granger e o de Wald HAC
mantêm o achado, porque o segundo corrige o erro-padrão, não a integração. O
de Toda-Yamamoto o apaga, e é ele o veredito. E o apaga nas \emph{duas}
direções: 0,45 no sentido Norte--Sul e 0,25 no sentido Sul--Norte.}

\pergunta{28. Como você sabe que aquela precedência era espúria e não real?}
\resposta{Por três passos, e nenhum deles foi ``a amostra é pequena''. A régua
estava errada para o material, o método robusto à integração apaga o resultado
nas duas direções, e --- o mais limpo --- a precedência reaparece onde não há
mecanismo que a explique: em três dos quatro placebos direcionais, incluindo o
pasto do Norte ``antecedendo'' o pasto do \emph{próprio Sul}, com o mesmo
p-valor do resultado original. Uma precedência que aparece em qualquer par de
séries suaves não é mecanismo, é defeito de instrumento.}

\pergunta{29. Como você controla multiplicidade?}
\resposta{Grades exploratórias são reportadas sob controle da taxa de falsas
descobertas, declarando quantos testes compõem a família, e o tamanho da
família altera o veredito, coisa que eu registro explicitamente: a célula
principal da Frente 3b sobrevive numa família de 96 e não sobrevive numa de 144.
Já a grade de dezesseis do teto de oferta \emph{não} recebe correção, porque são
dezesseis leituras da mesma hipótese pré-declarada, e não uma família de
hipóteses distintas. Aplicar FDR ali puniria a checagem em vez da pescaria.}

\pergunta{30. Você usa censo em vez de amostra. Como fica o teste de hipótese?}
\resposta{Perde o sentido usual, e essa é a decisão D23. Com $n$ igual à
população, qualquer desvio ínfimo do modelo nulo gera estatística
arbitrariamente grande --- a magnitude do $\Delta$BIC mede o tamanho de $n$, e
não a força da evidência. Nos módulos que rodam sobre o censo eu busco robustez
na \emph{estabilidade entre recortes}, com o mesmo padrão reaparecendo em
mesorregiões, AMCs e no pixel.}

\pergunta{31. Se o $\Delta$BIC não vale, o que sustenta as duas populações de
pastagem?}
\resposta{Evidências que não dependem de ajuste algum. Separados apenas pela
origem anterior do pixel, os eventos se ordenam sozinhos: pasto vindo de lavoura
volta a ser lavoura com mediana de 5 anos, e pasto vindo de vegetação natural
dura 13. Nas conversões de pasto com até três anos, 45\% já tinham sido lavoura
antes; nas de trinta anos ou mais, 1\%. O contraste de 45\% contra 1\% é o que o
dado estabelece com firmeza.}

\pergunta{32. A distribuição de idades não tem dois picos. Como você fala em
duas populações?}
\resposta{Porque a segunda é cerca de cinco vezes mais dispersa que a primeira,
com desvio de 7,5 anos contra 1,6, e espalha a mesma massa por uma faixa muito
maior --- ela produz um ombro, não uma segunda espiga. E eu não a estabeleço por
inspeção visual: o que a estabelece é que uma população única, ajustada sozinha,
devolve média de 12 anos e erra as duas pontas.}

\pergunta{33. A sua variável de esgotamento chegava a $-84{,}9$ numa fração que
deveria viver entre 0 e 1. Como isso passou?}
\resposta{Passou porque a variável entrava padronizada, e um z-score não é
robusto a nada: um único valor oitenta vezes fora de escala redefine a unidade
de medida da coluna inteira. Eram 14\% do painel, vindos de unidades cujo
estoque inicial era próximo de zero e que \emph{ganharam} estoque por oscilação
do classificador. O resultado foi o coeficiente achatado contra zero, e o nulo
que eu teria publicado. A D29 fixa a regra: toda variável com domínio declarado
é verificada contra ele antes de entrar padronizada.}

\pergunta{34. E como sei que a correção não foi escolhida por conveniência?}
\resposta{Porque não adotei um tratamento: rodei dezesseis combinações --- quatro
tratamentos, com e sem ponderação, em duas amostras --- e o coeficiente é
negativo nas dezesseis. Os dois tratamentos principais têm defeitos
\emph{diferentes} e chegam ao mesmo lugar: o de domínio descarta 46 unidades, e
o de piso em zero cria um acúmulo artificial em zero. É a concordância entre
defeitos independentes que carrega a leitura, não qualquer um deles.}

\pergunta{35. Um coeficiente de $-1{,}14$ na substituição local significa
substituição mais que completa. Isso não é suspeito?}
\resposta{É, e eu o digo assim. Ele não descreve substituição mais completa:
descreve mais de um hectare de pasto perdido por hectare de lavoura ganho, o que
por definição excede a substituição e inclui algum fluxo adicional não
identificado neste desenho. E ele vem da régua da união, que é o \emph{teto} do
intervalo da D26, não a medida certa. O coeficiente da régua estrita é
$-0{,}52$.}

\pergunta{36. O seu $\beta = +2{,}76$ do fluxo contra estoque tem $R^2$ de 0,43.
Não é o seu resultado mais forte?}
\resposta{Não é resultado nenhum. Fluxo é, por construção, taxa vezes estoque, e
segue daí sem nenhum dado que, se a taxa for aproximadamente constante, o fluxo
acompanha o estoque. Aquele número é a identidade se reapresentando, e eu o
reporto justamente para que ele não passe por achado. O conteúdo empírico está
na outra equação, a do comportamento da taxa.}

\pergunta{37. Acrescentar um termo quadrático não testaria não linearidade?}
\resposta{Testei, e ele sai nulo ($p = 0{,}92$). Isso diz que a relação entre
fluxo e estoque é a linha reta que a identidade prevê, sem curvatura própria ---
é uma afirmação sobre o estoque, e não sobre o comportamento da taxa.}

\pergunta{38. A sua variável de esgotamento e a sua taxa compartilham o mesmo
estoque. Isso não é circular?}
\resposta{Elas compartilham, sim --- no denominador de uma e no numerador com
sinal trocado da outra --- e vale registrar de que lado corre essa aritmética: o
acoplamento empurra o coeficiente para \emph{cima}. Achar sinal negativo é achar
contra a montagem, e não a favor dela. Ainda assim, o que eu afirmo é uma
regularidade compatível com atrito de oferta, e não que a dificuldade do
remanescente seja a causa da queda.}

\chapter{Sobre a interpretação e o alcance}

\pergunta{39. Então você refutou o iLUC?}
\resposta{Não, e isso é o que mais se perde na citação de terceiros. O canal que
testei é estreito: deslocamento intraestadual por adjacência, com defasagem
curta. O que posso dizer é que a assinatura testada não aparece nos dados
examinados, com o poder disponível. Ficam expressamente fora do alcance um
deslocamento de longo alcance que pule os vizinhos, um de defasagem longa, um
efeito temporal moderado, e qualquer canal que atravesse a divisa do estado.}

\pergunta{40. Se o Sul não empurra o Norte, o que explica eles marcharem
juntos?}
\resposta{A hipótese que testo é a de que ambos respondam ao mesmo estímulo
externo, com o câmbio como candidato mais forte. Mas ela não decorre de o
empurrão não ter aparecido --- descartar uma explicação não instala a seguinte.
Ela foi testada por conta própria, e o veredito é corroborante e não
estabelecido: o coeficiente tem a direção prevista e toda a bateria de
especificidade sai como deveria, mas o $p$ apropriado ao desenho não cruza os
5\%.}

\pergunta{41. E por que a aptidão deixou de ser a explicação?}
\resposta{Porque exógena não é o mesmo que isolada. A aptidão correlaciona-se
$-0{,}44$ com a latitude, e a latitude organiza em Goiás quase tudo o que varia
de sul para norte. Posta a latitude na mesma regressão, a aptidão perde 62\% da
magnitude e o seu $p$ de permutação vai de 0,13 para 0,47, e mesmo na régua
otimista sai de 0,026 para 0,30. A latitude quase não se move. Por isso o
resultado passou a ser enunciado como gradiente Sul--Norte, e a aptidão é a
régua com que ele foi medido, não o canal identificado.}

\pergunta{42. Na especificação com as três exposições, a latitude também deixa
de ser significativa. Isso não desfaz a sua decisão?}
\resposta{Não, porque a decisão se apoia em a aptidão \emph{perder}
significância diante do confundidor puro, e não em a latitude \emph{ganhá-la}. A
especificação com as três mostra a mesma coisa por outro lado: três gradientes
correlacionados e cerca de 38 realizações do choque não se separam entre si. E é
esse o veredito que eu reporto.}

\pergunta{43. O que impediria de dizer que o Sul goiano já fez a transição
florestal?}
\resposta{A definição. Numa das formulações, há transição quando a queda da
cobertura cessa \emph{e} a recuperação começa; noutra, a sequência tem quatro
estágios e o terceiro é a estabilização, separado do quarto, que é o
reflorestamento. O Sul goiano tem estoque reduzido, taxa cadente e estabilização
a partir de 2019, e \emph{não tem} regeneração líquida. Ele exibe o terceiro
estágio e não o quarto. A estabilização é esgotamento: perde-se pouco porque
restou pouco a perder.}

\pergunta{44. Cinco anos de estabilização não é pouco para essa leitura?}
\resposta{É, e eu o digo. Nenhuma das duas formulações fixa quanto tempo separa
a estabilização do reflorestamento, e sequência apenas inacabada e sequência
truncada produzem, nesse prazo, a mesma série. O que eu afirmo é a observação, e
não o desfecho: a estabilização registrada é esgotamento, e não autoriza, por
si, a leitura otimista que a fase tardia costuma sugerir.}

\pergunta{45. Você fala em ``teto de oferta'', mas 83\% da freada está numa
parcela residual. Isso não enfraquece o argumento?}
\resposta{Enfraqueceria se o argumento dependesse do tamanho da parcela de
estoque, e ele não depende. A parcela residual reúne tudo o que não é o tamanho
do estoque --- propensão a converter, atrito de acesso, proteção, troca de fonte
de terra --- e explicitamente \emph{não é uma medida de demanda}. O argumento se
monta por eliminação e por regularidade: a taxa do Sul caiu com a demanda no
pico, não foi a proteção integral, a demanda de terra continuou existindo e foi
atendida por pasto já aberto, e a taxa cai medidamente à medida que a unidade se
esgota.}

\pergunta{46. E se o que barrou a fronteira do Sul foi a lei, e não a falta de
terra?}
\resposta{Essa é a alternativa mais séria, e eu a separo em duas. A
\emph{proteção integral} eu elimino com dado: ela cobre menos de 3\% em qualquer
região, 76\% da que existe hoje já existia em 1985, e dos 6,56 Mha ainda
expostos 6,35 estão fora dela. Uma variável praticamente constante não data uma
inflexão, e ela não está à frente da fronteira. Já a Reserva Legal e a APP são
outra coisa: incidem dentro de cada propriedade, ganharam fiscalização com o CAR
a partir de 2012, e tendem a morder mais onde o remanescente já é escasso. No
Sul esgotado, ``acabou a terra convertível'' e ``o que restou é legalmente
inconversível'' produzem a mesma série. Separá-las exige cadastro ambiental
integrado pixel a pixel, que é a primeira aresta do meu cronograma.}

\pergunta{47. Os 6,56 Mha são terra disponível para a agricultura?}
\resposta{Não, e o vocabulário aqui é deliberado. O termo ``convertível'' mede
\emph{exposição}, e não disponibilidade: é o Cerrado que ainda pode ser perdido.
E é um teto, porque sem cadastro pixel a pixel o estoque inclui Reserva Legal e
APP que na prática não podem ser convertidas. A terra realmente exposta é menor,
nunca maior.}

\chapter{Sobre a literatura e a contribuição}

\pergunta{48. Qual é a contribuição original do trabalho?}
\resposta{Três coisas, em ordem de segurança. A primeira é submeter a teste
formal um canal que a literatura consultada não examinou: o deslocamento
intraestadual por adjacência --- a literatura de iLUC liga municípios distantes
centenas de quilômetros, e o canal de vizinhança dentro do mesmo estado não
aparece nas fontes lidas. A segunda é o desenho comparativo: Goiás reúne os dois
regimes de fronteira do Cerrado sob o mesmo classificador, a mesma malha de
território constante e as mesmas decisões metodológicas, o que separa Sul e
Norte não pode ser atribuído à régua, porque a régua é uma só. A terceira, mais
modesta e de método, é o centro de massa pesado por \emph{fluxo} de transição,
para o qual não localizei precedente.}

\pergunta{49. Só 38 referências não é pouco para uma dissertação?}
\resposta{É uma escolha declarada. A regra que adotei é que só entra nas
referências obra lida integralmente; o que o levantamento localizou e ainda não
foi lido está na §2.10, com forma bibliográfica completa e o que se espera de
cada uma. Uma bibliografia curta e conferida é mais defensável que uma longa e
citada de segunda mão. A revisão de literatura é, aliás, a frente mais extensa
do meu cronograma até a defesa, e ela parte exatamente dessa lista.}

\pergunta{50. Você diverge de Arima e colegas, que encontram o efeito na
Amazônia. Como explica?}
\resposta{A divergência delimita o que resta comparar, e não estabelece
contradição. Diferem a região, o desfecho medido e, sobretudo, a forma da
ligação espacial: o modelo deles liga municípios distantes centenas de
quilômetros, e o canal que eu testo é o de adjacência. E vale registrar que o
meu nulo \emph{não} é previsão cumprida: havia razão teórica para esperar o
efeito, e o precedente mais próximo o encontrou. Verificar se alguma literatura
já isolou esse canal é a frente prioritária da minha revisão --- dela depende a
leitura do resultado, não a sua obtenção.}

\chapter{As perguntas que podem doer}

Três perguntas exigem preparo específico, porque a resposta errada custa caro e
a resposta certa vale muito.

\section{``Então o seu trabalho não achou nada?''}

\begin{nabanca}
``Ele achou quatro coisas, e duas delas são nulos. O padrão está estabelecido e
sobrevive a pixel, malha, classe e barra de erro. O mecanismo local está
estabelecido: duas lógicas convivem em todo o estado e a região explica 0,5\% da
idade da pastagem. O empurrão entre regiões foi procurado em trinta e seis
recortes e não apareceu, e isso muda a política recomendada, porque se a
conversão no Norte fosse efeito dominó da lavoura do Sul, conter a lavoura
conteria a fronteira. E a freada do Sul não é mercado fraco: é compatível com
esgotamento, e eu tenho o mecanismo medido. Um nulo com poder declarado e
alternativa medida é resultado; o que não seria resultado é um silêncio.''
\end{nabanca}

\section{``Você usou inteligência artificial. O trabalho é seu?''}

Esta pergunta vai aparecer, porque o próprio texto a convida ao declarar o uso.
A resposta tem três partes e vale ensaiá-la inteira.

\begin{nabanca}[Primeiro, o escopo, com precisão]
``O escopo do uso foi a \emph{implementação}: escrita e depuração de código de
coleta e análise, execução de análises exploratórias, produção de figuras e
visualizações, e auditorias de consistência entre o que os dados dizem e o que
os documentos afirmam. A formulação das perguntas de pesquisa, as decisões
metodológicas e a interpretação dos resultados permanecem minhas. E essa
distinção não é retórica: as trinta e uma decisões numeradas são escolhas de
desenho, tomadas diante de alternativas com custos diferentes, e é sobre elas
que recai a responsabilidade autoral.''
\end{nabanca}

\begin{nabanca}[Segundo, a contrapartida de verificação]
``Um instrumento dessa natureza cobra contrapartida, e eu a organizei em quatro
práticas: o registro numerado das decisões, que permite reconstruir por que cada
escolha foi feita e o que ela substituiu; a rastreabilidade obrigatória de todo
número exibido até o arquivo que o produziu; a revisão cruzada entre assistentes
de fornecedores distintos, com verificação manual dos números por mim; e o
registro das autocorreções.''
\end{nabanca}

\begin{nabanca}[Terceiro, e é o que fecha: a limitação que eu mesmo medi]
``A revisão cruzada mostrou-se necessária porque essas revisões acertam
diagnósticos estruturais e ao mesmo tempo produzem números que \emph{parecem}
verificados e não são. Eu encontrei casos concretos disso e eles estão
documentados. É a razão de toda auditoria aqui ter escopo declarado, e de
nenhuma verificação automática do repositório conferir número contra fonte ---
isso é trabalho de leitura humana.''
\end{nabanca}

\begin{armadilha}[O que não dizer]
Não minimize (``usei só um pouco'') nem se desculpe. E não deixe a resposta
terminar na defesa: termine no que o arranjo \emph{produziu} --- 58 rotinas
reprodutíveis sobre dados públicos, um registro de decisões, sete autocorreções
documentadas e uma visualização navegável. Poucos trabalhos de mestrado chegam à
banca podendo mostrar o filtro, e não só o que passou por ele.
\end{armadilha}

\section{``Por que a sua hipótese inicial estava errada?''}

\begin{nabanca}
``Ela não estava errada; ela foi testada. A investigação partiu da hipótese
clássica da pastagem como estágio intermediário da ocupação, com a lavoura do
Sul empurrando a pecuária para o Norte. Parte dela se confirmou --- a pastagem
\emph{é} o elo intermediário de 58\% da entrada em agricultura. A parte mais
tentadora, o deslocamento causal de uma região sobre a outra, não se sustentou
quando submetida a um teste que podia falhar. O que emergiu no lugar organiza o
texto: duas dinâmicas paralelas, movidas pelo mesmo motor ao longo do gradiente
Sul--Norte e limitadas por um teto de oferta. A trajetória da pesquisa foi o
teste dessa hipótese, e é isso que eu apresento.''
\end{nabanca}
